\section*{Capítulo \ref{ch:g10:cells-common-unit} --- La célula: la unidad común de la vida}

\begin{solution}{exo:g10:cells-common-unit:1}
Todo ser vivo está hecho de una o de varias \omterm{def:g10:cells-common-unit:cell}{células}; la \omterm{def:g10:cells-common-unit:cell}{célula} es la
unidad básica de estructura y de función; toda \omterm{def:g10:cells-common-unit:cell}{célula} procede de una
\omterm{def:g10:cells-common-unit:cell}{célula} preexistente por división.
\end{solution}

\begin{solution}{exo:g10:cells-common-unit:2}
\qty{50}{\micro m}; \qty{2.5}{\micro m}; \qty{0.007}{mm}.
\end{solution}

\begin{solution}{exo:g10:cells-common-unit:3}
$36/600 = \qty{0.06}{mm} = \qty{60}{\micro m}$.
\end{solution}

\begin{solution}{exo:g10:cells-common-unit:4}
Solo vegetal: \omterm{def:g10:cells-common-unit:organelle}{pared celular}, \omterm{def:g10:cells-common-unit:organelle}{cloroplastos}, gran \omterm{def:g10:cells-common-unit:organelle}{vacuola} central. En
toda \omterm{def:g10:cells-common-unit:cell}{célula}: una \omterm{def:g10:cells-common-unit:cell}{membrana plasmática} (y también \omterm{def:g10:cells-common-unit:cell}{citoplasma} y
\omterm{def:g10:cells-common-unit:organelle}{ribosomas}).
\end{solution}

\begin{solution}{exo:g10:cells-common-unit:5}
Una \omterm{def:g10:cells-common-unit:prokaryote}{célula eucariota} guarda su ADN dentro de un \omterm{def:g10:cells-common-unit:organelle}{núcleo} y tiene
\omterm{def:g10:cells-common-unit:organelle}{orgánulos} rodeados de membrana (una \omterm{def:g10:cells-common-unit:cell}{célula} de mejilla, una levadura);
una \omterm{def:g10:cells-common-unit:prokaryote}{célula procariota} no tiene ni lo uno ni lo otro, su ADN está libre
en el \omterm{def:g10:cells-common-unit:cell}{citoplasma} y mide alrededor de \qty{1}{\micro m} (una bacteria
como \emph{E. coli}).
\end{solution}

\begin{solution}{exo:g10:cells-common-unit:6}
Longitud real $6 \times 5/15 = \qty{2}{\micro m}$. \omterm{def:g10:cells-common-unit:magnification}{Aumento}: la barra
mide \qty{15}{mm} para \qty{5}{\micro m} $= \qty{0.005}{mm}$, luego
$15/0.005 = 3000$: $\times 3000$.
\end{solution}

\begin{solution}{exo:g10:cells-common-unit:7}
Un \omterm{def:g10:cells-common-unit:organelle}{ribosoma} (\qty{25}{nm}) es diez veces menor que la \omterm{def:g10:cells-common-unit:magnification}{resolución} del
microscopio óptico (\qty{0.2}{\micro m} $= \qty{200}{nm}$): dos
\omterm{def:g10:cells-common-unit:organelle}{ribosomas} vecinos, o un \omterm{def:g10:cells-common-unit:organelle}{ribosoma} y su entorno, no pueden distinguirse.
Más \omterm{def:g10:cells-common-unit:magnification}{aumento} solo agranda el borrón.
\end{solution}

\begin{solution}{exo:g10:cells-common-unit:8}
Los granos son de almidón, guardado dentro de las \omterm{def:g10:cells-common-unit:cell}{células} en forma de
granos sólidos (el almidón es insoluble). Las reservas de una planta se
guardan dentro de sus \omterm{def:g10:cells-common-unit:cell}{células}, y no en algún espacio entre ellas.
\end{solution}

\begin{solution}{exo:g10:cells-common-unit:9}
La escama de cebolla crece bajo tierra, a oscuras, y acumula reservas:
allí los \omterm{def:g10:cells-common-unit:organelle}{cloroplastos} no servirían de nada. La hoja de musgo vive a la
luz y fabrica su alimento por fotosíntesis, así que sus \omterm{def:g10:cells-common-unit:cell}{células} están
repletas de \omterm{def:g10:cells-common-unit:organelle}{cloroplastos}. Mismo plan, distinto equipamiento.
\end{solution}

\begin{solution}{exo:g10:cells-common-unit:10}
Lado \qty{1}{\micro m}: superficie \qty{6}{\micro m^2}, volumen
\qty{1}{\micro m^3}, cociente 6 por micrómetro. Lado \qty{20}{\micro m}:
superficie \qty{2400}{\micro m^2}, volumen \qty{8000}{\micro m^3},
cociente 0,3. La bacteria tiene veinte veces más membrana por unidad de
contenido y vive solo de la difusión; la \omterm{def:g10:cells-common-unit:cell}{célula} mayor necesita
transporte interno y \omterm{def:g10:cells-common-unit:organelle}{orgánulos}.
\end{solution}

\begin{solution}{exo:g10:cells-common-unit:11}
\qty{2}{h} son 6 divisiones: $2^6 = 64$ \omterm{def:g10:cells-common-unit:cell}{células}. \qty{4}{h} son 12:
$2^{12} = 4096$.
\end{solution}

\begin{solution}{exo:g10:cells-common-unit:12}
Tiene membrana y \omterm{def:g10:cells-common-unit:cell}{citoplasma}, y toma y cede sustancias, así que encaja
en casi toda la definición, pero no puede dividirse ni renovar sus
\omterm{def:g10:chemistry-of-life:families}{proteínas}: es una \omterm{def:g10:cells-common-unit:cell}{célula} que ha renunciado a parte del programa a
cambio de sitio para la hemoglobina, y la fabrican y reponen \omterm{def:g10:cells-common-unit:cell}{células}
completas. La definición describe el caso general; el glóbulo rojo es
una forma especializada y terminal de él.
\end{solution}

\begin{solution}{exo:g10:cells-common-unit:13}
El corcho es tejido muerto: Hooke vio las paredes de celulosa de
\omterm{def:g10:cells-common-unit:cell}{células} cuyo contenido vivo había desaparecido --- cajas vacías. Un
\omterm{def:g10:cells-common-unit:organelle}{núcleo} solo existe en el \omterm{def:g10:cells-common-unit:cell}{citoplasma} de una \omterm{def:g10:cells-common-unit:cell}{célula} viva y, en cualquier
caso, la \omterm{def:g10:cells-common-unit:magnification}{resolución} de su microscopio y la falta de colorantes lo
habrían ocultado.
\end{solution}

\begin{solution}{exo:g10:cells-common-unit:14}
Un virus no tiene \omterm{def:g10:cells-common-unit:cell}{citoplasma}, ni membrana fabricada por él mismo, ni
metabolismo, y no puede dividirse: no cumple la definición de \omterm{def:g10:cells-common-unit:cell}{célula}.
Según la teoría celular no está vivo por su cuenta; es una partícula
que toma prestada la maquinaria de una \omterm{def:g10:cells-common-unit:cell}{célula} para hacer copias de sí
misma. Llamarlo o no ``vivo'' es cuestión de definición; que no es una
\omterm{def:g10:cells-common-unit:cell}{célula}, no lo es.
\end{solution}

\begin{solution}{exo:g10:cells-common-unit:15}
Volumen $\pi \times (25\,\unit{\micro m})^2 \times 30000\,\unit{\micro m}
\approx \num{5.9e7}\,\unit{\micro m^3}$, unas 7000 veces los
\qty{8000}{\micro m^3} del cubo. Al ser un cilindro delgado, ningún
punto de la fibra dista más de \qty{25}{\micro m} de la membrana, de
modo que los intercambios siguen siendo rápidos; y cada \omterm{def:g10:cells-common-unit:organelle}{núcleo} gobierna
solo su propio tramo de fibra, así que ninguno tiene que atender un
volumen mayor que el de una \omterm{def:g10:cells-common-unit:cell}{célula} corriente.
\end{solution}

\begin{solution}{pb:g10:cells-common-unit:1}
\textbf{1.} $1.8/9 = \qty{0.2}{mm} = \qty{200}{\micro m}$.

\textbf{2.} $1.8/4 = \qty{0.45}{mm}$; caben unas dos \omterm{def:g10:cells-common-unit:cell}{células}.

\textbf{3.} $28/0.2 = 140$: el dibujo es $\times 140$.

\textbf{4.} $4/140 \approx \qty{0.029}{mm}$, unos \qty{30}{\micro m}.

\textbf{5.} \omterm{def:g10:cells-common-unit:organelle}{Pared celular}, \omterm{def:g10:cells-common-unit:cell}{membrana plasmática} (pegada a la pared),
\omterm{def:g10:cells-common-unit:cell}{citoplasma}, \omterm{def:g10:cells-common-unit:organelle}{núcleo}, \omterm{def:g10:cells-common-unit:organelle}{vacuola}. Presentes pero invisibles a $\times 400$:
los \omterm{def:g10:cells-common-unit:organelle}{ribosomas} (y las \omterm{def:g10:cells-common-unit:organelle}{mitocondrias} están en el límite).

\textbf{6.} Vienen de la superficie de un revestimiento de varias
\omterm{def:g10:cells-common-unit:cell}{células} de grosor sometido a roce constante: las \omterm{def:g10:cells-common-unit:cell}{células} más externas
quedan aplanadas y se desprenden, y por eso las recoge un bastoncillo
de algodón.

\textbf{7.} \qty{0.24}{mm} está justo por encima del límite de
\qty{0.1}{mm} a simple vista: una mota apenas visible. A $\times 100$
ocupa $0.24/1.8 \approx 13\%$ del campo, y a $\times 400$ más de la
mitad; ningún objetivo disponible da exactamente un cuarto.

\textbf{8.} Bacterias. El microscopio electrónico mostraría su pared,
su membrana, sus \omterm{def:g10:cells-common-unit:organelle}{ribosomas} y su zona de ADN, y la ausencia de \omterm{def:g10:cells-common-unit:organelle}{núcleo}.

\textbf{9.} La \omterm{def:g10:cells-common-unit:cell}{célula} de mejilla hace un solo trabajo (proteger una
superficie) en un cuerpo donde otras \omterm{def:g10:cells-common-unit:cell}{células} la alimentan, la oxigenan
y la coordinan; el paramecio no tiene a nadie más y debe alimentarse,
moverse, percibir y dividirse como una sola \omterm{def:g10:cells-common-unit:cell}{célula}.

\textbf{10.} El exterior es ahora más salado que el \omterm{def:g10:cells-common-unit:cell}{citoplasma}, así que
el agua sale de la \omterm{def:g10:cells-common-unit:cell}{célula} a través de la membrana; la \omterm{def:g10:cells-common-unit:cell}{célula} pierde
volumen y se arruga.

\textbf{11.} $20^3 = \qty{8000}{\micro m^3} = \num{8e-6}\,\unit{mm^3}$.

\textbf{12.} $6 \times 20^2 = \qty{2400}{\micro m^2}$; cociente
$2400/8000 = 0.3$ por micrómetro.

\textbf{13.} Superficie \qty{6}{\micro m^2}, volumen \qty{1}{\micro m^3},
cociente 6: veinte veces mayor.

\textbf{14.} $\qty{8000}{\micro m^3} = \num{8e-9}\,\unit{cm^3}$, es
decir unos $\num{8e-9}\,\unit{g} = \qty{8}{ng}$.

\textbf{15.} $\pi \times 3.5^2 \times 2 \approx \qty{77}{\micro m^3}$,
alrededor de una centésima parte del cubo.

\textbf{16.} $\qty{60}{L} = \num{6e16}\,\unit{\micro m^3}$; dividido
entre 8000: unos $\num{7.5e12}$ \omterm{def:g10:cells-common-unit:cell}{células}.

\textbf{17.} Dos tercios de eso: unos $\num{5e12}$.

\textbf{18.} $\num{2.5e13} \times 77 \approx \num{1.9e15}\,\unit{\micro
m^3} \approx \qty{1.9}{L}$. \omterm{def:g10:cells-common-unit:cell}{Células} de la sangre: $0.45 \times 4.5 \approx
\qty{2.0}{L}$ --- coherente.

\textbf{19.} Unos $\num{5e12}$ \omterm{def:g10:cells-common-unit:cell}{células} grandes más $\num{2.5e13}$
glóbulos rojos: alrededor de $\num{3e13}$, unas decenas de billones.
Los tamaños celulares abarcan un factor de cien y la mezcla de clases
varía, así que todo recuento es una estimación con un factor de dos o
tres de incertidumbre.

\textbf{20.} Unos treinta billones ($\num{3e13}$) de \omterm{def:g10:cells-common-unit:cell}{células}, todas
producidas por divisiones celulares sucesivas a partir de una única
\omterm{def:g10:cells-common-unit:cell}{célula} inicial, el óvulo fecundado.
\end{solution}
